\section{Motivace vývojářů}

Jedním z klíčových prvků pro celkový úspěch projektu je motivace vývojářů. Studie \cite{developer-motivation, motivation-agile-teams} potvrzují, že motivovaní programátoři odvádějí lepší práci, což se pozitivně projevuje jak například na produktivitě vývojového týmu, tak i kvalitě dodávaného produktu nebo na dodržování rozpočtů. Z tohoto důvodu se firmy snaží hledat různé možnosti, jak motivovat své zaměstance.

Jeden z možných způsobů dělení motivace je na vnitřní a vnější. Mezi vnitřní motivátory můžeme u vývojářů například zařadit různorodost práce, řešení technických výzev, autonomii a prostor pro rozhodování nebo ztotožnění se zadaným úkolem. Vnější motivátory mohou naopak zahrnovat finanční odměny a benefity, vhodné pracovní podmínky, jistotu zaměstnání nebo rovnováhu mezi osobním a pracovním životem. Dle výsledků studií \cite{developer-motivation, motivation-agile-teams} jsou pro typického softwarového vývojáře důležitější vnitřní motivátory.

\begin{description}
    \item[Gamifikace] Gamifikace představuje využívání herních prvků a mechanik v neherních oblastech. Jejím cílem je zvýšit motivaci, aktivní zapojení uživatelů a angažovanost prostřednictvím technik známých z her. Mezi takové mechaniky je možné například zařadit získávání bodů nebo odznaků, sestavování žebříčků nebo plnění výzev. \cite{gamification}
\end{description}

Motivaci vývojářů lze tedy podpořit i gamifikačními prvky, které je možné zařadit jak mezi vnitřní, tak i mezi vnější motivátory. Příkladem gamifikace zaměřené na vnitřní motivaci by mohly být například výzvy, které budou vývojáře motivovat k učení se novým technologiím a dovednostem. Mezi vnější motivátory by naopak mohlo být zařazené již zmiňované získávání bodů a odznaků nebo sestavování žebříčků.

Konkrétním příkladem systému využívající gamifikační prvky by mohl být portál Stack Overflow, který vývojáři využívají pro pokládání technických dotazů. Za jejich zodpovězení mohou následně ostatní uživatelé získávat reputační body a odznaky, které se poté často zobrazují vedle jejich jména. Motivací vývojářů přispívajících na Stack Overflow se zabýval průzkum \cite{gamification-so}, jehož se zúčastnilo 938 programátorů. Více než polovina respondentů (62 \%) uvedlo, že je motivuje získávání reputačních bodů, a pro 38 \% respondentů má motivační efekt i zisk odznaků.
